\section{RL 与 Context：能力为何来自闭环}

如果模型只在静态语料上学习，它能吸收已有模式，却很难直接知道某次行动是否在现实任务中成功。强化学习把模型放入可交互环境，让它尝试行动、观察结果并根据 reward 更新策略。这里的 \term{RL} 是 Reinforcement Learning，中文通常译为强化学习；\term{environment} 是模型可以采取行动并收到状态或结果反馈的任务世界。两者把“更多 token”扩展为“更多经过验证的经验”。

\begin{figure}[H]
\centering
\includegraphics[width=0.9\textwidth]{images/00-19-45-how-rl-works.jpg}
\caption{强化学习闭环：环境给出状态，策略采取行动，反馈再推动策略更新。}
\figsource{Stanford Online 官方视频 00:19:45，`images/00-19-45-how-rl-works.jpg`。}
\end{figure}

\begin{knowledgebox}{读图：把三类对象分开}
图中的 agent 或 policy 是做决策的模型，environment 是会响应行动的外部系统，reward 是对结果的压缩反馈。三者不能互相替代：更大的模型不等于更好的环境，更丰富的环境不等于 reward 正确，reward 可计算也不等于它代表人类真正想要的结果。系统能力来自闭环质量，而不是某个组件的孤立规模。
\end{knowledgebox}

一个简化的 RL 目标可以写成
\[
J(\theta)=\mathbb{E}_{\tau\sim\pi_\theta}\left[\sum_{t=0}^{T}\gamma^t r_t\right].
\]
其中，$\pi_\theta$ 表示参数为 $\theta$ 的策略，$\tau$ 表示策略与环境交互形成的一条 trajectory，$r_t$ 是第 $t$ 步 reward，$\gamma$ 是对远期回报的折扣因子，$T$ 是任务终止时刻。公式没有回答 reward 应该怎样设计，也没有保证训练一定稳定；它只把“让高回报轨迹更可能出现”写成了优化对象。

\begin{warningbox}{课堂提示：经验 scaling 不等于无限 scaling}
讲者反复把 scaling laws 描述成 empirical observations。训练规模扩大后性能曾经持续提升，只能说明在已观察区间内存在规律。RL 也可能因环境质量、验证器漏洞、探索不足、分布偏移或目标错配而平台化。工程规划可以使用曲线外推，但必须同时准备失效条件和替代实验。
\end{warningbox}

\subsection{Basic AI Scaling Recipe}

本节把现代模型能力的生产过程压成四个节点：compute 支撑 research，research 产出 models，models 通过 inference 接触用户与任务，使用过程再产生 context 与反馈。这个图的重点不在箭头数量，而在闭环方向。训练不是一次性把知识封装进权重；部署后的环境、工具调用和成功失败记录会决定下一轮研究与训练能够看见什么。

\begin{figure}[H]
\centering
\includegraphics[width=0.76\textwidth]{images/00-22-34-basic-ai-scaling-recipe.jpg}
\caption{Basic AI Scaling Recipe：compute、research、models、inference 与 context feedback 的循环。}
\figsource{Stanford Online 官方视频 00:22:34，`images/00-22-34-basic-ai-scaling-recipe.jpg`。}
\end{figure}

\begin{importantbox}{读图：闭环中最容易被忽略的箭头}
从 compute 到 model 的路径容易量化，GPU 数量、训练 FLOPs 和模型规模都有数字。从 inference 回到 context feedback 的路径更难：需要知道用户在什么工作流里使用模型、何时失败、失败是否可验证、数据是否有权回流、反馈能否代表目标人群。讲者把 context 放在首位，是因为这一箭头经常决定后续 compute 有没有高价值用途。
\end{importantbox}

\subsection{Context 不只是更长的 prompt}

这里的 context 范围远大于对话窗口。它包括任务说明、私有文档、代码库、工具权限、执行环境、历史决策、用户行为和可验证结果。一个模型即使拥有相同权重，在能否访问终端、测试、部署日志与组织知识方面不同，也会表现成完全不同的系统。Frontier Systems 因此把 context 看作一种基础设施资产：它需要采集、授权、清洗、连接和持续维护。

\begin{figure}[H]
\centering
\includegraphics[width=0.9\textwidth]{images/00-24-04-context-environments-flywheel.jpg}
\caption{Context/Environments flywheel：接入工作流，观察失败，改进环境与模型，再扩大使用。}
\figsource{Stanford Online 官方视频 00:24:04，`images/00-24-04-context-environments-flywheel.jpg`。}
\end{figure}

\begin{knowledgebox}{读图：三个问题决定 flywheel 是否成立}
第一，系统接触的是高价值真实任务，还是只在演示数据上循环；第二，成功与失败能否被可靠验证，还是只能依赖模糊满意度；第三，数据回流是否得到用户授权并保持安全边界。只有使用增长、反馈质量与模型改进互相促进时，flywheel 才成立。单纯增加调用量可能只会放大噪声、隐私风险和错误自动化。
\end{knowledgebox}

\begin{table}[H]
\centering
\begin{tabular}{L{0.18\textwidth}L{0.32\textwidth}L{0.4\textwidth}}
\toprule
术语 & 它解决的问题 & 本讲中的含义 \\
\midrule
Prompt context & 当前请求需要什么信息 & 对话、文档与即时约束 \\
Tool context & 模型能做哪些动作 & API、终端、搜索、数据库与权限 \\
Environment & 动作在哪里执行并产生结果 & 代码仓库、浏览器、实验室或业务系统 \\
Feedback & 怎样判断行动效果 & 测试、用户修改、收入、延迟或安全事件 \\
Flywheel & 改进为何能持续积累 & 使用产生高质量反馈，反馈又改善模型与产品 \\
\bottomrule
\end{tabular}
\caption{术语表：把泛化的 context 拆成可设计、可审计的组件。}
\end{table}

\subsection{本章小结}

RL 提供了策略—环境—反馈的学习闭环，Basic Scaling Recipe 把这个闭环放进 compute、research、model 与 inference 的全栈过程。Context 的关键不在长度，而在任务、工具、环境和验证信号的质量。也正因为高价值 context 稀缺，模型公司与应用公司之间会围绕入口、数据和工作流产生激烈竞争。

